\documentclass[aps,preprint]{revtex4}%
\usepackage{amsfonts}
\usepackage{amsmath}
\usepackage{amssymb}
\usepackage{graphicx}%
\setcounter{MaxMatrixCols}{30}
%TCIDATA{OutputFilter=latex2.dll}
%TCIDATA{Version=4.10.0.2363}
%TCIDATA{CSTFile=revtex4.cst}
%TCIDATA{Created=Monday, June 24, 2002 09:35:27}
%TCIDATA{LastRevised=Thursday, September 11, 2003 11:27:57}
%TCIDATA{<META NAME="GraphicsSave" CONTENT="32">}
%TCIDATA{<META NAME="DocumentShell" CONTENT="Articles\SW\REVTeX 4">}
\newtheorem{theorem}{Theorem}
\newtheorem{acknowledgement}[theorem]{Acknowledgement}
\newtheorem{algorithm}[theorem]{Algorithm}
\newtheorem{axiom}[theorem]{Axiom}
\newtheorem{claim}[theorem]{Claim}
\newtheorem{conclusion}[theorem]{Conclusion}
\newtheorem{condition}[theorem]{Condition}
\newtheorem{conjecture}[theorem]{Conjecture}
\newtheorem{corollary}[theorem]{Corollary}
\newtheorem{criterion}[theorem]{Criterion}
\newtheorem{definition}[theorem]{Definition}
\newtheorem{example}[theorem]{Example}
\newtheorem{exercise}[theorem]{Exercise}
\newtheorem{lemma}[theorem]{Lemma}
\newtheorem{notation}[theorem]{Notation}
\newtheorem{problem}[theorem]{Problem}
\newtheorem{proposition}[theorem]{Proposition}
\newtheorem{remark}[theorem]{Remark}
\newtheorem{solution}[theorem]{Solution}
\newtheorem{summary}[theorem]{Summary}
\newenvironment{proof}[1][Proof]{\noindent\textbf{#1.} }{\ \rule{0.5em}{0.5em}}
\begin{document}
\preprint{ }
\title[TDVP-HJ ]{Hamilton - Jacobi Theory and the Time Dependent Variational Principle }
\author{Brian Weiner}
\affiliation{Pennsylvania State University}
\date{06/24/02}
\startpage{1}
\maketitle
\tableofcontents


\section{Introduction}

The Time Dependent Variational Principle has been a cornerstone in the
development of exact, i.e. the Time Dependent Schr\"{o}dinger Equation (TDSE),
and approximate quantum equations of motion derived from an action principle.
First introduced by Dirac \cite{dirac}, it was further developed by McLachlan
and Ball \cite{mac&ball} and then by Kramer and Saraceno \cite{kram&sar}. Two
basic approaches have been utilized in the derivation; one can either consider
pure states to be described in terms of unnormalized vectors which leads to an
unconstrained variational problem or in terms of normalized vectors which
leads to a constrained variational problem. Kramer and Saraceno used the
unconstrained approach and after the implementation of the boundary conditions
of fixed end points, i.e. the initial and final state vectors of the quantum
evolution are held fixed, obtained an equation for the time dependent state
vector that is a stationary point of the action. However, even in the case
when there are no restrictions on the state vectors, this equation is not the
TDSE and solutions to it are not solutions to the TDSE! In order to obtain
solutions of the TDSE they multiply solutions of this equation by a specific
time dependent factor, a complex phase factor, which can in fact be obtained
as a solution of an independent equation. Unfortunately the introduction of
this phase factor produces a state vector path that does not have fixed
endpoints in contradiction to the assumptions invoked to obtain equations for
the path where the action is stationary.

In this presentation I show what action principle and associated boundary
conditions actually do give the TDSE and at the same time take into account
the often ignored fact that a physical pure quantum state is not described by
a single vector in a Hilbert space but a line in that space.

The proper description of the action principle that generates the TDSE is
crucial in the

\begin{itemize}
\item examination of alternative forms of this equation such as in Time
Dependent Density Functional Theory

\item construction of approximations to the TDSE equation

\item discussions of causality
\end{itemize}


\end{document}